% GENERATED FILE. Edit canonical lessons, then run npm run generate:books.
% canonical-source-hash: fnv1a64:dabcb9b76e2a5a85
% canonical-lessons: ML-C453-vottu-ceyyuka, ML-C453-pintunaykkuka, ML-C453-etirppu, ML-C453-codyam-ceyyuka, ML-C453-pratikarikkuka

\chapter{For and Against}
\label{ch:ml-for-and-against}
\hlchaptermodality{453}

\begin{chapteropening}
\textbf{By the end of this chapter:} \emph{I can say to vote, to support, an objection, to question and to respond.}

Complete the last lesson of chapter 453: I can say to vote, to support, an objection, to question and to respond.
\end{chapteropening}

\section[\texorpdfstring{\ml{വോട്ട്} \ml{ചെയ്യുക} (vōṭṭŭ ceyyuka)}{വോട്ട് ചെയ്യുക (vōṭṭŭ ceyyuka)}]{\ml{വോട്ട്} \ml{ചെയ്യുക} (vōṭṭŭ ceyyuka) — to vote}

\label{lesson:ML-C453-vottu-ceyyuka}



\begin{quote}
Before the new one: say the Malayalam for an approach, then the Malayalam for cooperation.
\end{quote}

\subsection*{You'll want to know: \ml{വോട്ട്} \ml{ചെയ്യുക}}
\textbf{\ml{വോട്ട്} \ml{ചെയ്യുക}} — \emph{vōṭṭŭ ceyyuka} — \textquotedblleft{}to vote\textquotedblright{}.

\begin{quote}
\ml{ഞാൻ} \ml{ഇന്നലെ} \ml{ആദ്യമായി} \ml{വോട്ട്} \ml{ചെയ്തു}.
\end{quote}

\begin{quote}
\emph{ñān innale ādyamāyi vōṭṭŭ ceytu.}
\end{quote}

\begin{quote}
Yesterday I voted for the first time.
\end{quote}

Past: \ml{വോട്ട്} \ml{ചെയ്തു}, \emph{vōṭṭŭ ceytu}.

For someone: -\ml{ന്} \ml{വോട്ട്} \ml{ചെയ്യുക}.

\begin{cousinweb}
\textbf{\ml{വോട്ട്} \ml{ചെയ്യുക}} (\emph{vōṭṭŭ ceyyuka}) is a phrase verb built on \ml{വോട്ട്}, which was borrowed from English \emph{vote}, with \ml{ചെയ്യുക}, \textquotedblleft{}to do\textquotedblright{}.
\end{cousinweb}

\begin{grammarlens}[title={What you've built}]
One more word for saying what happens if.
\end{grammarlens}

\subsection*{Your turn}
\begin{itemize}
\raggedright
  \item \emph{Say it:} \emph{vōṭṭŭ ceyyuka}
  \item \emph{Say it:} \emph{vōṭṭŭ ceyyuka}, once more
  \item \emph{Say it:} \emph{ñān innale ādyamāyi vōṭṭŭ ceytu.} — Yesterday I voted for the first time.
  \item \emph{Recall:} say the Malayalam for an approach, then the Malayalam for cooperation, then say \emph{vōṭṭŭ ceyyuka} again
\end{itemize}

\begin{tcolorbox}[breakable,colback=teal!4,colframe=teal!35!black,title={Before you move on}]
What does \ml{വോട്ട്} \ml{ചെയ്യുക} mean? (\textquotedblleft{}To vote\textquotedblright{}.) Say it once more.
\end{tcolorbox}
\section[\texorpdfstring{\ml{പിന്തുണയ്ക്കുക} (pintuṇaykkuka)}{പിന്തുണയ്ക്കുക (pintuṇaykkuka)}]{\ml{പിന്തുണയ്ക്കുക} (pintuṇaykkuka) — to stand by; to support}

\label{lesson:ML-C453-pintunaykkuka}



\begin{quote}
Before the new one: say the Malayalam for cooperation, then the Malayalam for to vote.
\end{quote}

\subsection*{You'll want to know: \ml{പിന്തുണയ്ക്കുക}}
\textbf{\ml{പിന്തുണയ്ക്കുക}} — \emph{pintuṇaykkuka} — \textquotedblleft{}to stand by; to support\textquotedblright{}.

\begin{quote}
\ml{നാട്ടുകാർ} \ml{അവളുടെ} \ml{പദ്ധതിയെ} \ml{പിന്തുണച്ചു}.
\end{quote}

\begin{quote}
\emph{nāṭṭukār avaḷuṭe paddhatiye pintuṇaccu.}
\end{quote}

\begin{quote}
The local people supported her plan.
\end{quote}

Past: \ml{പിന്തുണച്ചു}, \emph{pintuṇaccu}.

Its past is \ml{പിന്തുണച്ചു}. What is supported takes -\ml{നെ}.

\begin{cousinweb}
\textbf{\ml{പിന്തുണയ്ക്കുക}} (\emph{pintuṇaykkuka}) is native Dravidian, the verb of \ml{പിന്തുണ}: \ml{പിൻ}, \textquotedblleft{}behind\textquotedblright{} (DEDR 4205), and \ml{തുണ}, \textquotedblleft{}help\textquotedblright{} (DEDR 3308), whose verb \emph{tuṇekka}, \textquotedblleft{}to help\textquotedblright{}, DEDR also lists.
\end{cousinweb}

\begin{grammarlens}[title={What you've built}]
One more word for saying what happens if.
\end{grammarlens}

\subsection*{Your turn}
\begin{itemize}
\raggedright
  \item \emph{Say it:} \emph{pintuṇaykkuka}
  \item \emph{Say it:} \emph{pintuṇaykkuka}, once more
  \item \emph{Say it:} \emph{nāṭṭukār avaḷuṭe paddhatiye pintuṇaccu.} — The local people supported her plan.
  \item \emph{Recall:} say the Malayalam for cooperation, then the Malayalam for to vote, then say \emph{pintuṇaykkuka} again
\end{itemize}

\begin{tcolorbox}[breakable,colback=teal!4,colframe=teal!35!black,title={Before you move on}]
What does \ml{പിന്തുണയ്ക്കുക} mean? (\textquotedblleft{}To stand by; to support\textquotedblright{}.) Say it once more.
\end{tcolorbox}
\section[\texorpdfstring{\ml{എതിർപ്പ്} (etirppŭ)}{എതിർപ്പ് (etirppŭ)}]{\ml{എതിർപ്പ്} (etirppŭ) — an objection; resistance}

\label{lesson:ML-C453-etirppu}



\begin{quote}
Before the new one: say the Malayalam for to vote, then the Malayalam for to stand by.
\end{quote}

\subsection*{You'll want to know: \ml{എതിർപ്പ്}}
\textbf{\ml{എതിർപ്പ്}} — \emph{etirppŭ} — \textquotedblleft{}an objection; resistance\textquotedblright{}.

\begin{quote}
\ml{പുതിയ} \ml{നികുതിക്കെതിരെ} \ml{വലിയ} \ml{എതിർപ്പ്} \ml{ഉണ്ടായി}.
\end{quote}

\begin{quote}
\emph{putiya nikutikketire valiya etirppŭ uṇṭāyi.}
\end{quote}

\begin{quote}
There was strong resistance to the new tax.
\end{quote}

The noun of \ml{എതിർക്കുക}, to oppose.

\begin{cousinweb}
\textbf{\ml{എതിർപ്പ്}} (\emph{etirppŭ}) is native Dravidian, the noun of \ml{എതിർക്കുക}, from \emph{etir}, \textquotedblleft{}that which is opposite, in front\textquotedblright{} (DEDR 795).
\end{cousinweb}

\begin{grammarlens}[title={What you've built}]
One more word for saying what happens if.
\end{grammarlens}

\subsection*{Your turn}
\begin{itemize}
\raggedright
  \item \emph{Say it:} \emph{etirppŭ}
  \item \emph{Say it:} \emph{etirppŭ}, once more
  \item \emph{Say it:} \emph{putiya nikutikketire valiya etirppŭ uṇṭāyi.} — There was strong resistance to the new tax.
  \item \emph{Recall:} say the Malayalam for to vote, then the Malayalam for to stand by, then say \emph{etirppŭ} again
\end{itemize}

\begin{tcolorbox}[breakable,colback=teal!4,colframe=teal!35!black,title={Before you move on}]
What does \ml{എതിർപ്പ്} mean? (\textquotedblleft{}An objection; resistance\textquotedblright{}.) Say it once more.
\end{tcolorbox}
\section[\texorpdfstring{\ml{ചോദ്യം} \ml{ചെയ്യുക} (cōdyaṁ ceyyuka)}{ചോദ്യം ചെയ്യുക (cōdyaṁ ceyyuka)}]{\ml{ചോദ്യം} \ml{ചെയ്യുക} (cōdyaṁ ceyyuka) — to question; to challenge}

\label{lesson:ML-C453-codyam-ceyyuka}



\begin{quote}
Before the new one: say the Malayalam for to stand by, then the Malayalam for an objection.
\end{quote}

\subsection*{You'll want to know: \ml{ചോദ്യം} \ml{ചെയ്യുക}}
\textbf{\ml{ചോദ്യം} \ml{ചെയ്യുക}} — \emph{cōdyaṁ ceyyuka} — \textquotedblleft{}to question; to challenge\textquotedblright{}.

\begin{quote}
\ml{പോലീസ്} \ml{അവനെ} \ml{മണിക്കൂറുകളോളം} \ml{ചോദ്യം} \ml{ചെയ്തു}.
\end{quote}

\begin{quote}
\emph{pōlīsŭ avane maṇikkūṟukaḷōḷaṁ cōdyaṁ ceytu.}
\end{quote}

\begin{quote}
The police questioned him for hours.
\end{quote}

Past: \ml{ചോദ്യം} \ml{ചെയ്തു}, \emph{cōdyaṁ ceytu}.

Also to challenge a decision: \ml{തീരുമാനം} \ml{ചോദ്യം} \ml{ചെയ്യുക}.

\begin{cousinweb}
\textbf{\ml{ചോദ്യം} \ml{ചെയ്യുക}} (\emph{cōdyaṁ ceyyuka}) is a phrase verb built on \ml{ചോദ്യം}, \textquotedblleft{}a question\textquotedblright{}, from Sanskrit \emph{codya}, \textquotedblleft{}to be urged on\textquotedblright{}, on the root \emph{cud}, \textquotedblleft{}to impel\textquotedblright{} (Monier-Williams), with \ml{ചെയ്യുക}, \textquotedblleft{}to do\textquotedblright{}.
\end{cousinweb}

\begin{grammarlens}[title={What you've built}]
One more word for saying what happens if.
\end{grammarlens}

\subsection*{Your turn}
\begin{itemize}
\raggedright
  \item \emph{Say it:} \emph{cōdyaṁ ceyyuka}
  \item \emph{Say it:} \emph{cōdyaṁ ceyyuka}, once more
  \item \emph{Say it:} \emph{pōlīsŭ avane maṇikkūṟukaḷōḷaṁ cōdyaṁ ceytu.} — The police questioned him for hours.
  \item \emph{Recall:} say the Malayalam for to stand by, then the Malayalam for an objection, then say \emph{cōdyaṁ ceyyuka} again
\end{itemize}

\begin{tcolorbox}[breakable,colback=teal!4,colframe=teal!35!black,title={Before you move on}]
What does \ml{ചോദ്യം} \ml{ചെയ്യുക} mean? (\textquotedblleft{}To question; to challenge\textquotedblright{}.) Say it once more.
\end{tcolorbox}
\section[\texorpdfstring{\ml{പ്രതികരിക്കുക} (pratikarikkuka)}{പ്രതികരിക്കുക (pratikarikkuka)}]{\ml{പ്രതികരിക്കുക} (pratikarikkuka) — to respond; to react}

\label{lesson:ML-C453-pratikarikkuka}



\begin{quote}
Before the new one: say the Malayalam for an objection, then the Malayalam for to question.
\end{quote}

\subsection*{You'll want to know: \ml{പ്രതികരിക്കുക}}
\textbf{\ml{പ്രതികരിക്കുക}} — \emph{pratikarikkuka} — \textquotedblleft{}to respond; to react\textquotedblright{}.

\begin{quote}
\ml{ആ} \ml{വാർത്തയോട്} \ml{മന്ത്രി} \ml{ഉടനെ} \ml{പ്രതികരിച്ചു}.
\end{quote}

\begin{quote}
\emph{ā vārttayōṭŭ mantri uṭane pratikariccu.}
\end{quote}

\begin{quote}
The minister responded to that news at once.
\end{quote}

Past: \ml{പ്രതികരിച്ചു}, \emph{pratikariccu}.

To something: -\ml{ഓട്} \ml{പ്രതികരിക്കുക}. The noun is \ml{പ്രതികരണം}, a response.

\begin{cousinweb}
\textbf{\ml{പ്രതികരിക്കുക}} (\emph{pratikarikkuka}) is built from Sanskrit \emph{prati-kṛ}, \textquotedblleft{}to act back\textquotedblright{}, whose noun \emph{pratikāra} means \textquotedblleft{}requital, counteraction\textquotedblright{} (Monier-Williams), with the Malayalam verb ending -\ml{ിക്കുക}.
\end{cousinweb}

\begin{grammarlens}[title={What you've built}]
One more word for saying what happens if.
\end{grammarlens}

\subsection*{Your turn}
\begin{itemize}
\raggedright
  \item \emph{Say it:} \emph{pratikarikkuka}
  \item \emph{Say it:} \emph{pratikarikkuka}, once more
  \item \emph{Say it:} \emph{ā vārttayōṭŭ mantri uṭane pratikariccu.} — The minister responded to that news at once.
  \item \emph{Recall:} say the Malayalam for an objection, then the Malayalam for to question, then say \emph{pratikarikkuka} again
\end{itemize}

\begin{tcolorbox}[breakable,colback=teal!4,colframe=teal!35!black,title={Before you move on}]
What does \ml{പ്രതികരിക്കുക} mean? (\textquotedblleft{}To respond; to react\textquotedblright{}.) Say it once more.
\end{tcolorbox}
